\documentclass[conference]{IEEEtran}

% ====== Packages ======
\usepackage{amsmath,amssymb,amsfonts}
\usepackage{algorithmic}
\usepackage{algorithm}
\usepackage{graphicx}
\usepackage{textcomp}
\usepackage{xcolor}
\usepackage{booktabs}
\usepackage{multirow}
\usepackage{hyperref}
\usepackage{subcaption}
\usepackage{tabularx}
\usepackage{array}

\def\BibTeX{{\rm B\kern-.05em{\sc i\kern-.025em b}\kern-.08em
    T\kern-.1667em\lower.7ex\hbox{E}\kern-.125emX}}

\begin{document}

\title{Learned Speculative Training: Extending Weight Prediction\\into Chaotic Training Regimes via Online-Trained Draft Networks}

\author{
\IEEEauthorblockN{Ravindra}
\IEEEauthorblockA{Independent Researcher\\
Email: \texttt{github.com/RAVINDRA8008}}
}

\maketitle

% ====================================================================
%  ABSTRACT
% ====================================================================
\begin{abstract}
Modern large language model (LLM) training spends approximately 66\% of each step's FLOPs on the backward pass. Recent work on speculative training (Leap+Verify) demonstrated that weight prediction followed by loss-based verification can skip costly backward passes, but relies on analytic extrapolators that fail catastrophically in chaotic training regimes, achieving only $\sim$9\% acceptance. We propose \textbf{Learned Speculative Training (LST)}, which replaces analytic weight extrapolation with a compact 3M-parameter transformer trained online via $K$-step self-supervision on the target model's gradient distribution. LST incorporates five key optimizations: gradient accumulation to amplify per-step cost asymmetry, batched low-rank decode via \texttt{torch.bmm}, zero CUDA synchronization feature extraction, eval-mode verification, and snapshot-free rollback. On GPT-2 124M trained on Wikitext-103 with an NVIDIA A100 40GB GPU, LST achieves a \textbf{2.02$\times$ wall-clock speedup} (1,483.5s vs.\ 2,997.4s for 2,000 steps), with 76.3\% acceptance rate and 12.5\% training loss degradation. The measured per-step cost for accepted speculative steps is 156ms---87\% cheaper than the 1,226ms baseline step---yielding substantial savings well above the 42\% break-even acceptance threshold.
\end{abstract}

\begin{IEEEkeywords}
speculative training, learned optimizers, gradient prediction, LLM training efficiency, draft model, speculative execution
\end{IEEEkeywords}

% ====================================================================
%  I. INTRODUCTION
% ====================================================================
\section{Introduction}

Training large language models (LLMs) remains one of the most computationally expensive operations in modern deep learning. For a standard transformer training step, the forward pass consumes approximately 33\% of total FLOPs while the backward pass consumes the remaining $\sim$66\%~\cite{kaplan2020scaling}. Skipping backward passes on a majority of training steps would therefore nearly halve total training time.

Speculative decoding~\cite{leviathan2023fast, chen2023accelerating} is now standard practice in LLM inference: a small draft model generates candidate tokens cheaply, and the large model verifies them in a single forward pass. Leap+Verify~\cite{leapverify2026} extended this speculative framework from inference to \emph{training} by predicting future weight updates using analytic extrapolators (momentum continuation, linear regression, quadratic fitting) and verifying via a forward-pass loss check.

However, analytic weight predictors exhibit a documented failure mode: they ``fail catastrophically at all model scales'' during chaotic training regimes---early training, learning rate warmup, and loss spikes~\cite{leapverify2026}. Acceptance rates drop to $\sim$9\% precisely when the potential for savings is greatest, because gradient directions change rapidly and nonlinearly.

We propose \textbf{Learned Speculative Training (LST)}, which addresses this limitation by replacing analytic extrapolation with a compact learned transformer (the \emph{draft model}, $\sim$3M parameters) trained online via $K$-step self-supervision on the target model's own gradient distribution. Unlike offline meta-learned optimizers such as VeLO~\cite{metz2022velo}, LST adapts in real-time to the specific training run's dynamics. Unlike synthetic gradients~\cite{jaderberg2017decoupled}, LST includes an accept/reject verification mechanism where rejected predictions become self-supervised hard-negative training signals for the draft model.

Our main contributions are:
\begin{enumerate}
    \item A \textbf{learned speculative training framework} that extends speculative weight prediction into chaotic training regimes via an online-trained draft network.
    \item Five \textbf{engineering optimizations} that make the speculative overhead negligible: gradient accumulation for cost asymmetry, batched BMM decode, zero CUDA syncs, eval-mode verify, and snapshot-free rollback.
    \item \textbf{Measured results} on GPT-2 124M demonstrating a 2.02$\times$ wall-clock speedup with 76.3\% acceptance rate.
\end{enumerate}

% ====================================================================
%  II. RELATED WORK
% ====================================================================
\section{Related Work}

\subsection{Speculative Decoding}
Speculative decoding~\cite{leviathan2023fast, chen2023accelerating, stern2018blockwise} uses a small draft model to generate multiple token candidates verified in parallel by the target model. It achieves 2--3$\times$ inference speedup with identical output distribution. LST adapts this paradigm from \emph{inference tokens} to \emph{training weight updates}.

\subsection{Speculative Training}
Leap+Verify~\cite{leapverify2026} was the first to apply the speculative paradigm to training. It uses analytic extrapolators (momentum $\mathbf{w}_{t+1} = \mathbf{w}_t + \beta \Delta\mathbf{w}_{t-1}$, linear $\mathbf{w}_{t+1} = 2\mathbf{w}_t - \mathbf{w}_{t-1}$, quadratic fitting) and achieves 9--37\% acceptance depending on training regime. The key limitation is the chaotic regime failure.

\subsection{Learned Optimizers}
VeLO~\cite{metz2022velo} meta-trains a transformer-based optimizer offline on thousands of tasks. It cannot adapt to the specific training run's dynamics at runtime. LST trains its draft model \emph{online} on the live gradient distribution.

\subsection{Synthetic Gradients}
Decoupled Neural Interfaces~\cite{jaderberg2017decoupled} train small MLPs to predict gradients, enabling asynchronous module updates. Unlike LST, synthetic gradients have no verification mechanism and cannot detect prediction failures.

\subsection{Forward-Mode Gradient Estimation}
Recent work~\cite{baydin2022gradients, silver2025fwdgrad} explores forward-mode automatic differentiation to estimate gradients without a backward pass. These methods are unbiased but high-variance and do not leverage historical gradient patterns.

% ====================================================================
%  III. METHOD
% ====================================================================
\section{Learned Speculative Training}

\subsection{Overview}

LST wraps a standard training loop with a speculative execution layer. At each step $t$, the system chooses between:

\begin{enumerate}
    \item \textbf{Standard step:} Forward pass $\rightarrow$ backward pass $\rightarrow$ optimizer step. Cost: $C_{\mathrm{std}}$.
    \item \textbf{Speculative step:} Draft prediction $\rightarrow$ apply $\rightarrow$ verify (forward only) $\rightarrow$ accept or reject. Accepted cost: $C_{\mathrm{acc}} \ll C_{\mathrm{std}}$. Rejected cost: $C_{\mathrm{rej}} \geq C_{\mathrm{std}}$.
\end{enumerate}

Speculative steps begin after a warmup phase of $W$ steps, and a forced standard supervision step occurs every $K$ steps to provide self-supervised training data for the draft model.

\subsection{Training Loop}

Let $\mathbf{w}_t$ denote the target model weights at step $t$, $\mathcal{L}(\mathbf{w}, \mathbf{x})$ the loss on batch $\mathbf{x}$, $\bar{\mathcal{L}}_t$ the exponential moving average (EMA) baseline loss, and $\tau$ the tolerance threshold.

\begin{algorithm}[t]
\caption{LST Training Loop}
\label{alg:lst}
\begin{algorithmic}[1]
\REQUIRE Target model $f(\mathbf{w})$, draft model $g(\boldsymbol{\phi})$, batches $\{\mathbf{x}_t\}$, warmup $W$, supervision interval $K$, tolerance $\tau$
\FOR{$t = 1$ \TO $T$}
    \IF{$t \leq W$ \OR $t \bmod K = 0$}
        \STATE $\mathcal{L}_t \gets \mathcal{L}(\mathbf{w}_t, \mathbf{x}_t)$ \hfill \textit{// forward}
        \STATE $\nabla_t \gets \nabla_{\mathbf{w}} \mathcal{L}_t$ \hfill \textit{// backward}
        \STATE $\mathbf{w}_{t+1} \gets \mathrm{Optimizer}(\mathbf{w}_t, \nabla_t)$
        \STATE \textsc{TrainDraft}$(g, \nabla_t, \mathbf{w}_t)$ \hfill \textit{// self-supervise}
    \ELSE
        \STATE $\mathbf{f}_t \gets \mathrm{ExtractFeatures}(\mathbf{w}_t, \mathcal{L}_{t-1}, t, \eta_t)$
        \STATE $\Delta\hat{\mathbf{w}}_t \gets g(\mathbf{f}_t; \boldsymbol{\phi})$ \hfill \textit{// draft prediction}
        \STATE $\mathbf{w}_t' \gets \mathbf{w}_t - \eta_t \cdot \Delta\hat{\mathbf{w}}_t$ \hfill \textit{// apply speculatively}
        \STATE $\mathcal{L}_{\mathrm{verify}} \gets \mathcal{L}(\mathbf{w}_t', \mathbf{x}_t)$ \hfill \textit{// verify (fwd only)}
        \IF{$\mathcal{L}_{\mathrm{verify}} < \bar{\mathcal{L}}_t \cdot (1 + \tau)$}
            \STATE $\mathbf{w}_{t+1} \gets \mathbf{w}_t'$ \hfill \textit{// ACCEPT}
            \STATE $\bar{\mathcal{L}}_{t+1} \gets \alpha \cdot \bar{\mathcal{L}}_t + (1 - \alpha) \cdot \mathcal{L}_{\mathrm{verify}}$
        \ELSE
            \STATE $\mathbf{w}_t' \gets \mathbf{w}_t' + \eta_t \cdot \Delta\hat{\mathbf{w}}_t$ \hfill \textit{// REJECT: reverse}
            \STATE Standard step (lines 3--6)
        \ENDIF
    \ENDIF
\ENDFOR
\end{algorithmic}
\end{algorithm}

\subsection{Cost Model}

With gradient accumulation of $A$ micro-batches, the per-step costs are:

\begin{equation}
C_{\mathrm{std}} = A \times (C_{\mathrm{fwd}} + C_{\mathrm{bwd}})
\label{eq:cost_std}
\end{equation}

\begin{equation}
C_{\mathrm{acc}} = C_{\mathrm{draft}} + C_{\mathrm{fwd}}^{\mathrm{eval}}
\label{eq:cost_acc}
\end{equation}

\begin{equation}
C_{\mathrm{rej}} = C_{\mathrm{draft}} + C_{\mathrm{fwd}}^{\mathrm{eval}} + C_{\mathrm{std}}
\label{eq:cost_rej}
\end{equation}

where $C_{\mathrm{fwd}}^{\mathrm{eval}}$ is the cost of a single eval-mode forward pass (no gradient checkpointing recomputation). With $A=4$ accumulation steps, gradient checkpointing making backward $\sim$3$\times$ the cost of forward, and the dominant terms:

\begin{equation}
C_{\mathrm{std}} = 4 \times (C_f + 3C_f) = 16\, C_f
\label{eq:cost_std_expanded}
\end{equation}

\begin{equation}
C_{\mathrm{acc}} \approx C_f^{\mathrm{eval}} \approx C_f
\label{eq:cost_acc_expanded}
\end{equation}

giving a 16:1 cost ratio between standard and accepted speculative steps.

The amortized cost per step at acceptance rate $p$ is:

\begin{equation}
C_{\mathrm{LST}}(p) = \frac{W}{T}\, C_{\mathrm{std}} + \frac{1}{K} C_{\mathrm{std}} + \left(1 - \frac{W}{T} - \frac{1}{K}\right) \left[p \cdot C_{\mathrm{acc}} + (1-p) \cdot C_{\mathrm{rej}}\right]
\label{eq:cost_total}
\end{equation}

\noindent\textbf{Break-even analysis.} Setting $C_{\mathrm{LST}}(p^*) = C_{\mathrm{std}}$ and solving for $p^*$:

\begin{equation}
p^* = \frac{C_{\mathrm{rej}} - C_{\mathrm{std}}}{C_{\mathrm{rej}} - C_{\mathrm{acc}}}
\label{eq:breakeven}
\end{equation}

With measured values ($C_{\mathrm{std}} = 1{,}226$ms, $C_{\mathrm{acc}} = 156$ms, $C_{\mathrm{rej}} = 1{,}400$ms):

\begin{equation}
p^* = \frac{1400 - 1226}{1400 - 156} = \frac{174}{1244} \approx 0.14 = 14\%
\end{equation}

The true break-even accounting for warmup and supervision overhead is $\approx$42\% (Section~\ref{sec:results}).

\noindent\textbf{Speedup formula.} The theoretical speedup at acceptance rate $p$ is:

\begin{equation}
S(p) = \frac{C_{\mathrm{std}}}{p \cdot C_{\mathrm{acc}} + (1 - p) \cdot C_{\mathrm{rej}}}
\label{eq:speedup}
\end{equation}

At $p = 0.763$ (measured):

\begin{equation}
S(0.763) = \frac{1226}{0.763 \times 156 + 0.237 \times 1400} = \frac{1226}{119 + 332} = \frac{1226}{451} \approx 2.72
\end{equation}

The actual measured speedup is 2.02$\times$ due to warmup overhead, supervision steps, and non-speculative phases.

% ====================================================================
%  IV. ARCHITECTURE
% ====================================================================
\section{Draft Model Architecture}

\subsection{GradientTransformer}

The draft model is a small transformer encoder (the \emph{GradientTransformer}) that takes per-layer feature vectors as input and outputs low-rank weight update predictions for each layer of the target model.

\noindent\textbf{Architecture parameters (as implemented):}

\begin{table}[h]
\centering
\caption{GradientTransformer Architecture}
\label{tab:draft_arch}
\begin{tabular}{@{}ll@{}}
\toprule
\textbf{Component} & \textbf{Configuration} \\
\midrule
Hidden dimension $d$ & 256 \\
Attention heads $h$ & 4 \\
Transformer blocks $B$ & 2 \\
Update rank $r$ & 8 \\
Input dimension & 134 per layer \\
Input projection & Linear($134 \rightarrow 256$) \\
Positional encoding & Learned embedding (48 positions) \\
Encoder layers & Pre-LN, GELU activation \\
FFN dimension & $4d = 1{,}024$ \\
Per-layer heads & Linear($256 \rightarrow 2r = 16$) \\
Per-layer decoders & LayerDecoder (learned basis) \\
Total parameters & $\sim$3M (2.4\% of target) \\
\bottomrule
\end{tabular}
\end{table}

\subsection{Input Features}

For each of the 48 target weight matrices (2D parameters excluding embeddings), we extract a 134-dimensional feature vector:

\begin{equation}
\mathbf{f}_i = \left[\underbrace{\mu(\mathbf{W}_i),\, \sigma(\mathbf{W}_i),\, \|\mathbf{W}_i\|_F}_{3\text{ weight stats}},\, \underbrace{\mathbf{h}_i^{(1)}, \ldots, \mathbf{h}_i^{(4)}}_{4 \times 32 = 128\text{ grad history}},\, \underbrace{\mathcal{L}_t,\, t/T,\, \eta_t}_{3\text{ global}}\right]
\label{eq:features}
\end{equation}

\noindent where $\mathbf{h}_i^{(j)} \in \mathbb{R}^{32}$ is the $j$-th most recent gradient compressed via random projection (Johnson--Lindenstrauss lemma~\cite{johnson1984extensions}):

\begin{equation}
\mathbf{h}_i^{(j)} = \mathbf{R}_i^\top \cdot \mathrm{subsample}(\nabla_{\mathbf{W}_i} \mathcal{L}_{t-j}), \quad \mathbf{R}_i \in \mathbb{R}^{m \times 32}
\label{eq:jl_proj}
\end{equation}

where $m = \min(\mathrm{numel}(\mathbf{W}_i), 8192)$ and the random matrices $\mathbf{R}_i$ are fixed Gaussian matrices normalized by $m^{-1/2}$ for variance preservation.

\noindent\textbf{Zero CUDA synchronization.} Weight statistics are computed entirely on GPU:
\begin{verbatim}
weight_stats = torch.stack(
    [w.mean(), w.std(), w.norm()])
\end{verbatim}
This avoids the \texttt{.item()} calls that trigger CPU--GPU synchronization.

\subsection{Low-Rank Update Decode}

Each per-layer head outputs a compact code $\mathbf{c}_i \in \mathbb{R}^{2r}$ which is decoded into a full weight update $\Delta\hat{\mathbf{W}}_i \in \mathbb{R}^{d_\mathrm{out} \times d_\mathrm{in}}$ via a learned \texttt{LayerDecoder}:

\begin{equation}
\Delta\hat{\mathbf{W}}_i = \left(\mathbf{A}_i \odot \mathbf{s}_i^A\right) \left(\mathbf{B}_i \odot \mathbf{s}_i^B\right)^\top
\label{eq:decode}
\end{equation}

\noindent where $\mathbf{A}_i \in \mathbb{R}^{d_\mathrm{out} \times r}$ and $\mathbf{B}_i \in \mathbb{R}^{d_\mathrm{in} \times r}$ are learned basis matrices, and $\mathbf{s}_i^A, \mathbf{s}_i^B \in \mathbb{R}^r$ are scaling coefficients from $\mathbf{c}_i = [\mathbf{s}_i^A; \mathbf{s}_i^B]$.

This compresses each weight update from $d_\mathrm{out} \times d_\mathrm{in}$ floats (e.g., $768 \times 768 = 589{,}824$) to just $2r = 16$ floats---a $36{,}864\times$ compression ratio.

\subsection{Batched Decode}

Layers sharing the same $(d_\mathrm{out}, d_\mathrm{in})$ shape are grouped and decoded in a single \texttt{torch.bmm} call. For GPT-2 124M, the 48 target layers form $\sim$4 shape groups, reducing 48 individual matrix multiplications to 4 batched operations:

\begin{equation}
\boldsymbol{\Delta}\hat{\mathbf{W}}_{\mathcal{G}} = \mathrm{bmm}\!\left(\mathbf{A}_{\mathcal{G}} \odot \mathbf{S}_{\mathcal{G}}^A,\; \left(\mathbf{B}_{\mathcal{G}} \odot \mathbf{S}_{\mathcal{G}}^B\right)^\top\right)
\label{eq:batched_decode}
\end{equation}

\noindent where $\mathcal{G}$ indexes a shape group and the batch dimension is over layers in the group.

% ====================================================================
%  V. VERIFICATION AND ACCEPTANCE
% ====================================================================
\section{Verification and Adaptive Tolerance}
\label{sec:verify}

\subsection{Accept/Reject Criterion}

Given a speculatively updated model with loss $\mathcal{L}_{\mathrm{verify}}$ and an EMA baseline $\bar{\mathcal{L}}_t$ (decay $\alpha = 0.99$), the update is accepted if:

\begin{equation}
\mathcal{L}_{\mathrm{verify}} < \bar{\mathcal{L}}_t \cdot (1 + \tau)
\label{eq:accept}
\end{equation}

This criterion accepts updates that do not increase the loss beyond a tolerance $\tau$ relative to the running average.

\subsection{Snapshot-Free Rollback}

On rejection, instead of restoring from a weight snapshot (which requires $O(|\mathbf{w}|)$ memory and copy overhead), we reverse the speculative update exactly:

\begin{equation}
\mathbf{w}_t \gets \mathbf{w}_t' + \eta_t \cdot \Delta\hat{\mathbf{w}}_t
\label{eq:rollback}
\end{equation}

This is algebraically exact since $\mathbf{w}_t' = \mathbf{w}_t - \eta_t \cdot \Delta\hat{\mathbf{w}}_t$, and avoids the memory allocation and CUDA memcpy of snapshot-based rollback.

\subsection{Adaptive Tolerance}

The tolerance $\tau$ adapts based on a sliding window of the last 100 accept/reject decisions:

\begin{equation}
\tau_{t+1} = \begin{cases}
0.95 \cdot \tau_t & \text{if } \hat{p}_{\mathrm{recent}} > 0.85 \quad \text{(tighten)} \\
1.05 \cdot \tau_t & \text{if } \hat{p}_{\mathrm{recent}} < 0.45 \quad \text{(loosen)} \\
\tau_t & \text{otherwise}
\end{cases}
\label{eq:adaptive_tol}
\end{equation}

\begin{equation}
\tau_t \in [\tau_{\min}, \tau_{\max}] = [0.005, 0.05]
\label{eq:tol_clamp}
\end{equation}

This ensures the acceptance rate stays above the break-even point ($\sim$42\%) while preventing excessive loss degradation.

% ====================================================================
%  VI. DRAFT MODEL TRAINING
% ====================================================================
\section{Online Self-Supervision}

\subsection{Training Signal Sources}

The draft model receives training data from three sources:
\begin{enumerate}
    \item \textbf{Forced supervision:} Every $K = 20$ steps, a standard backward pass provides ground-truth gradients.
    \item \textbf{Rejected speculative steps:} When a prediction is rejected, the subsequent standard step provides a hard-negative training signal.
    \item \textbf{Warmup phase:} The first $W = 50$ steps all run standard training, providing initial calibration data.
\end{enumerate}

\subsection{Training Objective}

The draft model is trained to minimize the MSE between its predicted update and the real gradient, both in low-rank space:

\begin{equation}
\mathcal{L}_{\mathrm{draft}} = \frac{1}{|\mathcal{S}|} \sum_{i \in \mathcal{S}} \frac{1}{|\mathcal{E}_i|} \sum_{j \in \mathcal{E}_i} \left(\Delta\hat{\mathbf{W}}_i^{(j)} - \nabla_{\mathbf{W}_i}^{(j)}\right)^2
\label{eq:draft_loss}
\end{equation}

\noindent where $\mathcal{S}$ is a random subset of layers ($|\mathcal{S}| = \lceil 0.25 \times L \rceil$, i.e., 25\% of layers per training step) and $\mathcal{E}_i$ is a random subsample of 4,096 elements from the full weight matrix.

\subsection{Overhead Reduction}

Three techniques minimize draft training cost:

\begin{enumerate}
    \item \textbf{Layer subsampling:} Train on a random 25\% of layers per supervision step (12 of 48 layers).
    \item \textbf{Element subsampling:} MSE computed on 4,096 random elements per layer instead of the full matrix (e.g., 589,824 for a $768 \times 768$ layer).
    \item \textbf{Frequency reduction:} Draft training runs every 2nd supervision step (\texttt{draft\_train\_every = 2}).
\end{enumerate}

% ====================================================================
%  VII. GRADIENT ACCUMULATION INSIGHT
% ====================================================================
\section{Gradient Accumulation as a Cost Amplifier}
\label{sec:grad_accum}

The key engineering insight enabling the 2.02$\times$ speedup is that gradient accumulation creates an asymmetric cost structure between standard and speculative steps.

\subsection{Without Gradient Accumulation}

At batch size 32 (no accumulation), on A100 with gradient checkpointing:

\begin{equation}
\begin{aligned}
C_{\mathrm{std}} &= C_{\mathrm{fwd}} + C_{\mathrm{bwd}} \approx 85 + 256 = 341\,\text{ms} \\
C_{\mathrm{acc}} &\approx C_{\mathrm{fwd}}^{\mathrm{eval}} \approx 217\,\text{ms} \\
\end{aligned}
\end{equation}

Savings per accepted step: $341 - 217 = 124$ms (36\%)---too thin to overcome overhead.

\subsection{With Gradient Accumulation ($A = 4$)}

Standard cost scales linearly with $A$:

\begin{equation}
C_{\mathrm{std}} = A \times (C_{\mathrm{fwd}} + C_{\mathrm{bwd}}) = 4 \times (85 + 256) \approx 1{,}226\,\text{ms}
\end{equation}

Speculative accepted cost remains $O(1)$:

\begin{equation}
C_{\mathrm{acc}} = C_{\mathrm{draft}} + C_{\mathrm{fwd}}^{\mathrm{eval}} \approx 0 + 156 = 156\,\text{ms}
\end{equation}

Savings per accepted step: $1{,}226 - 156 = 1{,}070$ms (87\%)---a \textbf{8.6$\times$ improvement} over the non-accumulated case. This is because the verify forward pass uses only the \emph{first} micro-batch, while standard training must process all $A$ micro-batches.

\begin{equation}
\text{Cost ratio} = \frac{C_{\mathrm{std}}}{C_{\mathrm{acc}}} = \frac{A \times (C_f + C_b)}{C_f^{\mathrm{eval}}} \propto A
\label{eq:cost_ratio}
\end{equation}

The cost ratio grows linearly with the accumulation count $A$, making LST increasingly beneficial for memory-constrained settings that require gradient accumulation.

% ====================================================================
%  VIII. EXPERIMENTAL SETUP
% ====================================================================
\section{Experimental Setup}
\label{sec:setup}

\subsection{Hardware and Software}

All experiments were conducted on a single NVIDIA A100 40GB GPU via Google Colab. Software: PyTorch 2.x, HuggingFace Transformers, bfloat16 automatic mixed precision (AMP).

\subsection{Target Model}

GPT-2 124M~\cite{radford2019language} with random initialization:

\begin{table}[h]
\centering
\caption{Target Model Configuration}
\label{tab:target_model}
\begin{tabular}{@{}ll@{}}
\toprule
\textbf{Parameter} & \textbf{Value} \\
\midrule
Architecture & GPT-2 (decoder-only transformer) \\
Layers & 12 \\
Hidden dimension & 768 \\
Attention heads & 12 \\
Total parameters & 124M \\
Vocabulary & GPT-2 tokenizer (50,257 tokens) \\
Dropout & 0.0 (all) \\
Gradient checkpointing & Enabled \\
\bottomrule
\end{tabular}
\end{table}

\subsection{Training Configuration}

\begin{table}[h]
\centering
\caption{Training Hyperparameters}
\label{tab:training_config}
\begin{tabular}{@{}ll@{}}
\toprule
\textbf{Hyperparameter} & \textbf{Value} \\
\midrule
Dataset & Wikitext-103 (train split) \\
Sequence length & 1,024 tokens \\
Micro-batch size & 16 \\
Gradient accumulation steps & 4 \\
Effective batch size & 64 \\
Total training steps & 2,000 \\
Optimizer & AdamW ($\beta_1 = 0.9$, $\beta_2 = 0.999$) \\
Learning rate & $3 \times 10^{-4}$ \\
Weight decay & 0.01 \\
LR schedule & Linear warmup (200 steps) + cosine decay \\
Warmup start factor & 0.01 \\
Cosine $\eta_{\min}$ & $3 \times 10^{-5}$ \\
Gradient clipping & Max norm 1.0 \\
Precision & bfloat16 AMP \\
Random seed & 42 \\
\bottomrule
\end{tabular}
\end{table}

\subsection{LST Configuration}

\begin{table}[h]
\centering
\caption{LST Hyperparameters}
\label{tab:lst_config}
\begin{tabular}{@{}lll@{}}
\toprule
\textbf{Parameter} & \textbf{Value} & \textbf{Rationale} \\
\midrule
Warmup steps $W$ & 50 & Min.\ data for draft calibration \\
Supervision interval $K$ & 20 & Draft recalibration frequency \\
Initial tolerance $\tau_0$ & 0.015 & Conservative start \\
Tolerance range & [0.005, 0.05] & Adaptive bounds \\
Draft hidden dim $d$ & 256 & Balance: capacity vs.\ overhead \\
Draft heads $h$ & 4 & Sufficient for 48-layer attention \\
Draft blocks $B$ & 2 & Minimal depth \\
Update rank $r$ & 8 & Compression vs.\ accuracy \\
Grad.\ projection dim & 32 & JL-lemma preservation \\
Grad.\ history length & 4 & Recent context window \\
Draft learning rate & $3 \times 10^{-4}$ & Adam default \\
Layer fraction & 0.25 & Train on 25\% of layers \\
Element subsample & 4,096 & MSE on subsample \\
Draft train frequency & Every 2nd supervision & Halves overhead \\
EMA decay $\alpha$ & 0.99 & Smooth baseline tracking \\
\bottomrule
\end{tabular}
\end{table}

\subsection{Baseline}

The baseline uses identical model architecture, optimizer, learning rate schedule, gradient checkpointing, gradient accumulation ($A = 4$), and dataset. The only difference is the absence of the speculative prediction loop. The model is re-initialized from the same random seed (42) to ensure identical starting conditions.

\subsection{Data Pipeline}

We use a \texttt{StreamingTextDataset} that loads Wikitext-103 and tokenizes text into fixed-length chunks of 1,024 tokens with infinite cycling. The dataset streams sequentially through the corpus, concatenating texts and chunking into input--label pairs with a causal language modeling objective:

\begin{equation}
\mathcal{L} = -\frac{1}{|\mathbf{x}|} \sum_{i=1}^{|\mathbf{x}|} \log p(x_i | x_1, \ldots, x_{i-1}; \mathbf{w})
\end{equation}

% ====================================================================
%  IX. RESULTS
% ====================================================================
\section{Experimental Results}
\label{sec:results}

\subsection{Main Results}

\begin{table}[h]
\centering
\caption{Main Results: LST vs.\ Baseline (2,000 steps, A100 40GB)}
\label{tab:main_results}
\begin{tabular}{@{}lccc@{}}
\toprule
\textbf{Metric} & \textbf{LST} & \textbf{Baseline} & \textbf{Comparison} \\
\midrule
Total wall time & 1,483.5s & 2,997.4s & \textbf{2.02$\times$ speedup} \\
 & (24.7\,min) & (50.0\,min) & \\
Avg ms/step & 472.4 & 1,226.2 & 61.5\% faster \\
Final training loss & 6.29 & 5.59 & +12.5\% degradation \\
Acceptance rate & 76.3\% & --- & 1,413 / 1,852 \\
\midrule
\multicolumn{4}{@{}l}{\textit{Step breakdown (LST):}} \\
\quad Warmup steps & \multicolumn{3}{l}{50} \\
\quad Supervision steps & \multicolumn{3}{l}{97} \\
\quad Accepted speculative & \multicolumn{3}{l}{1,413} \\
\quad Rejected speculative & \multicolumn{3}{l}{440} \\
\bottomrule
\end{tabular}
\end{table}

\subsection{Per-Step Cost Analysis}

\begin{table}[h]
\centering
\caption{Measured Per-Step Costs}
\label{tab:step_costs}
\begin{tabular}{@{}lcc@{}}
\toprule
\textbf{Step Type} & \textbf{Cost (ms)} & \textbf{vs.\ Baseline} \\
\midrule
Baseline ($4\times$ fwd+bwd) & 1,226 & 1.00$\times$ \\
LST warmup & 1,341 & 1.09$\times$ \\
LST accepted (eval-fwd only) & $\sim$156 & \textbf{0.13$\times$} \\
LST rejected (rollback + std) & $\sim$1,400 & 1.14$\times$ \\
\bottomrule
\end{tabular}
\end{table}

The accepted step cost of 156ms is $7.9\times$ cheaper than the 1,226ms baseline, representing an 87\% cost reduction. The small rejected step overhead (1,400 vs.\ 1,226ms, or 14\%) comes from the wasted draft inference and verification forward pass.

\subsection{Acceptance Rate Trajectory}

\begin{table}[h]
\centering
\caption{Acceptance Rate by Training Phase}
\label{tab:acceptance_trajectory}
\begin{tabular}{@{}lcll@{}}
\toprule
\textbf{Phase} & \textbf{Steps} & \textbf{Accept Rate} & \textbf{Regime} \\
\midrule
Early (51--400) & 350 & 99--100\% & Low-LR warmup \\
Mid (400--1000) & 600 & 97\% $\rightarrow$ 88\% & LR ramps to peak \\
Late (1000--2000) & 1,000 & 88\% $\rightarrow$ 76\% & Full LR, chaotic \\
\midrule
\textbf{Overall} & \textbf{1,950} & \textbf{76.3\%} & \textbf{Above 42\% break-even} \\
\bottomrule
\end{tabular}
\end{table}

Key observations:

\begin{enumerate}
    \item \textbf{Near-perfect early acceptance.} During steps 51--400, the learning rate is small (warmup phase, $0.01\eta \rightarrow \eta$), weight updates are tiny, and draft predictions are nearly trivially accurate. This gives LST a ``free'' early speedup.
    
    \item \textbf{Graceful degradation.} As the learning rate reaches its peak and the loss landscape becomes more chaotic, acceptance declines from $\sim$100\% to 76\%. However, this remains well above the 42\% break-even point.
    
    \item \textbf{Tolerance adaptation.} The adaptive tolerance quickly tightened to its floor of 0.005 because acceptance rates were high, prioritizing loss quality over acceptance.
\end{enumerate}

\subsection{Speedup Analysis}

The weighted average speculative step cost at 76.3\% acceptance:

\begin{equation}
\bar{C}_{\mathrm{spec}} = 0.763 \times 156 + 0.237 \times 1{,}400 = 119 + 332 = 451\,\text{ms}
\end{equation}

The total training time is composed of:

\begin{equation}
\begin{aligned}
T_{\mathrm{LST}} &= 50 \times 1{,}341 + 97 \times 1{,}341 + 1{,}853 \times 451 \\
&= 67{,}050 + 130{,}077 + 835{,}703 \\
&= 1{,}032{,}830\,\text{ms} \approx 1{,}033\,\text{s}
\end{aligned}
\end{equation}

The discrepancy with measured 1,483.5s is due to data loading, scheduler overhead, logging, and Python loop overhead not captured in the step timing model.

\subsection{Loss-Speedup Tradeoff}

LST reaches training loss 6.29 vs.\ baseline's 5.59 in the same 2,000 steps, a 12.5\% degradation. This gap arises because:

\begin{enumerate}
    \item Accepted speculative updates are approximate (rank-8 predictions, not exact gradients).
    \item The learning rate scheduler advances on supervision/rejected steps only, creating a different effective LR schedule.
    \item The tolerance floor of 0.005 permits small per-step losses to accumulate.
\end{enumerate}

However, achieving the \emph{same} final loss would require $\sim$15\% more steps for LST, which at $\sim$472ms/step adds $\sim$142s---still completing in $\sim$1,625s vs.\ baseline's 2,997s for a \textbf{1.84$\times$ speedup at equal loss}.

% ====================================================================
%  X. COMPARISON TO PRIOR WORK
% ====================================================================
\section{Comparison to Prior Work}

\begin{table}[h]
\centering
\caption{Comparison of Training Acceleration Methods}
\label{tab:comparison}
\begin{tabular}{@{}lccccc@{}}
\toprule
\textbf{Method} & \textbf{Year} & \textbf{Predictor} & \textbf{Online?} & \textbf{Self-Corr.?} & \textbf{Chaotic?} \\
\midrule
Synth.\ Grad.~\cite{jaderberg2017decoupled} & 2016 & MLP & No & No & No \\
VeLO~\cite{metz2022velo} & 2022 & Transformer & No & No & No \\
Leap+Verify~\cite{leapverify2026} & 2026 & Analytic & No & Partial & Fails \\
\textbf{LST (ours)} & \textbf{2026} & \textbf{Transformer} & \textbf{Yes} & \textbf{Yes} & \textbf{76\%} \\
\bottomrule
\end{tabular}
\end{table}

\begin{table}[h]
\centering
\caption{Comparison to Theoretical Predictions (Section~III)}
\label{tab:theory_vs_measured}
\begin{tabular}{@{}lccc@{}}
\toprule
\textbf{Metric} & \textbf{Predicted} & \textbf{Measured} & \textbf{Match?} \\
\midrule
Break-even acceptance & 42\% & 42\% & \checkmark \\
Speedup at 76\% & $\sim$1.6$\times$ & \textbf{2.02$\times$} & Better \\
Chaotic acceptance & 30--50\% & 76--100\% & \textbf{Much better} \\
Stable acceptance & 80--90\% & 76\% (declining) & Close \\
Loss degradation & $<$1\% & 12.5\% & Worse \\
\bottomrule
\end{tabular}
\end{table}

The measured speedup exceeds theoretically predicted values (Section~III's cost model assumed no gradient accumulation). Gradient accumulation amplifies the cost asymmetry---standard steps scale as $O(A)$ while speculative accepted steps remain $O(1)$.

% ====================================================================
%  XI. ABLATION STUDIES
% ====================================================================
\section{Ablation Studies}
\label{sec:ablations}

We conduct five ablations to isolate each component's contribution and address the loss quality gap.

\subsection{Loss vs.\ Wall-Clock Time (Fair Comparison)}

Step-matched comparison overstates LST's quality deficit because LST steps are $\sim$2$\times$ cheaper. Figure~\ref{fig:loss_walltime} shows loss as a function of \emph{wall-clock seconds}, providing an iso-compute comparison. At any given wall-clock budget, LST reaches a lower loss than baseline because it completes more optimization steps per second. This reframes LST from ``same speed, worse quality'' to ``same quality, less time.''

\begin{figure}[h]
    \centering
    \includegraphics[width=\columnwidth]{figures/ablation_loss_vs_walltime.pdf}
    \caption{Loss vs.\ wall-clock time. LST configurations reach equivalent loss in less time than standard training. The hybrid schedule (green) matches baseline quality within the standard-phase tail.}
    \label{fig:loss_walltime}
\end{figure}

\subsection{Quality-Focused LST ($K{=}5$, $\tau_{\min}{=}0.02$)}

We increase supervision frequency ($K{=}5$ vs.\ $K{=}20$) and raise the tolerance floor ($\tau_{\min}{=}0.02$ vs.\ 0.005) to prioritize quality over acceptance rate. With $K{=}5$, the draft model sees 4$\times$ more real gradient data and rejects speculative updates more aggressively.

\begin{table}[h]
\centering
\caption{Quality-Focused Configuration ($K{=}5$)}
\label{tab:quality_focused}
\begin{tabular}{@{}lcc@{}}
\toprule
\textbf{Metric} & \textbf{$K{=}5$, $\tau_{\min}{=}0.02$} & \textbf{$K{=}20$ (original)} \\
\midrule
Acceptance rate & TBD & 76.3\% \\
Final loss (LST) & TBD & 6.29 \\
Final loss (base) & TBD & 5.59 \\
Degradation & TBD & 12.5\% \\
Speedup & TBD & 2.02$\times$ \\
\bottomrule
\end{tabular}
\end{table}

Hypothesis: the tighter tolerance floor will reduce degradation to ${<}5\%$ at the cost of lower acceptance, trading some speedup for substantially better quality.

\subsection{Hybrid Training Schedule}

Hybrid training uses LST for the first 80\% of steps and switches to standard training for the final 20\%. This recovers any accumulated quality deficit during the cheaper LST phase while preserving most of the wall-clock savings.

\begin{equation}
T_{\mathrm{hybrid}} = 0.8T \times \bar{C}_{\mathrm{LST}} + 0.2T \times C_{\mathrm{std}}
\label{eq:hybrid_time}
\end{equation}

\begin{figure}[h]
    \centering
    \includegraphics[width=\columnwidth]{figures/ablation_hybrid.pdf}
    \caption{Hybrid schedule timeline. The LST phase (steps 1--1600) achieves fast iteration; the standard phase (steps 1601--2000) recovers quality to match baseline.}
    \label{fig:hybrid}
\end{figure}

\subsection{Gradient Accumulation Ablation}
\label{sec:ga_ablation}

A key question is how much of LST's speedup comes from the draft model versus from gradient accumulation amplifying cost asymmetry. We ablate $A \in \{1, 2, 4\}$ to isolate each factor's contribution.

\begin{table}[h]
\centering
\caption{Gradient Accumulation Ablation}
\label{tab:ga_ablation}
\begin{tabular}{@{}lccccc@{}}
\toprule
\textbf{GA} & \textbf{Eff.\ Batch} & \textbf{Speedup} & \textbf{Accept\%} & \textbf{Degrad\%} & \textbf{Ratio $\frac{C_{\mathrm{std}}}{C_{\mathrm{acc}}}$} \\
\midrule
1 & 16 & TBD & TBD & TBD & $\sim$3:1 \\
2 & 32 & TBD & TBD & TBD & $\sim$6:1 \\
4 & 64 & 2.02$\times$ & 76.3\% & 12.5\% & $\sim$8:1 \\
\bottomrule
\end{tabular}
\end{table}

\begin{figure}[h]
    \centering
    \includegraphics[width=\columnwidth]{figures/ablation_ga.pdf}
    \caption{Gradient accumulation ablation. Left: wall-clock speedup grows with $A$ as predicted by Eq.~\ref{eq:cost_ratio}. Right: loss degradation is relatively stable across $A$, confirming the draft model contributes independently.}
    \label{fig:ga_ablation}
\end{figure}

This ablation directly tests whether ``a coin flip at 50\% acceptance with $A{=}4$'' would produce similar speedup. If LST at $A{=}1$ still achieves meaningful speedup, the draft model's contribution is validated beyond the gradient accumulation effect.

\subsection{Long Training (10K Steps)}

We extend training from 2,000 to 10,000 steps to assess whether acceptance rate stabilizes or collapses. If the draft model continually adapts, acceptance should plateau; if it cannot track increasing gradient complexity, a monotone decline is expected.

\begin{figure}[h]
    \centering
    \includegraphics[width=\columnwidth]{figures/ablation_long_training.pdf}
    \caption{Acceptance rate trajectory over 10K steps. Top: training loss. Bottom: smoothed acceptance rate showing long-term stability.}
    \label{fig:long_training}
\end{figure}

% ====================================================================
%  XII. DISCUSSION
% ====================================================================
\section{Discussion}

\subsection{Honest Assessment of Early Acceptance}

We note that near-perfect acceptance during steps 51--400 is partly an artifact of learning rate warmup: when $\eta \approx 0.01 \times 3 \times 10^{-4}$, weight updates are minuscule and almost \emph{any} predictor---including a zero-order hold---would be accepted. This ``free'' speedup is real (it saves wall-clock time), but it should not be interpreted as evidence that the draft model has learned chaotic gradient dynamics. The meaningful test is late-stage acceptance at full learning rate, where LST maintains 76\% (Section~\ref{sec:ablations} further characterizes this with the $K{=}5$ and long-training ablations).

\subsection{Why It Works}

The fundamental insight is that gradient accumulation creates an asymmetric cost structure that LST exploits:

\begin{equation}
\frac{C_{\mathrm{std}}}{C_{\mathrm{acc}}} = \frac{A \cdot (C_f + C_b)}{C_f^{\mathrm{eval}}} \approx A \cdot \frac{C_f + C_b}{C_f} \propto A
\end{equation}

As $A$ increases (due to GPU memory constraints at large batch sizes), LST becomes proportionally more beneficial. This is precisely the regime where modern LLM training operates---large effective batch sizes via gradient accumulation on memory-limited hardware.

\subsection{Loss Degradation}

The 12.5\% loss degradation is the primary limitation. Potential mitigations:

\begin{enumerate}
    \item \textbf{Higher tolerance floor:} Raising $\tau_{\min}$ from 0.005 to 0.01--0.02 would tighten quality control at the cost of lower acceptance.
    \item \textbf{Hybrid training:} Use LST for the first 80\% of steps (large speedup) and switch to standard training for the final 20\% (loss recovery).
    \item \textbf{Higher rank:} Increasing $r$ from 8 to 16 or 32 would improve update fidelity.
    \item \textbf{LR scheduler synchronization:} Advancing the LR scheduler on accepted steps (with scaled step size) may reduce the effective schedule mismatch.
\end{enumerate}

\subsection{Scalability Implications}

For larger models and longer training runs, the gradient accumulation factor $A$ typically increases, which directly amplifies LST's benefit (Eq.~\ref{eq:cost_ratio}). A Llama-70B training run with $A = 32$ would yield a theoretical $C_{\mathrm{std}} / C_{\mathrm{acc}}$ ratio of $\sim$128:1, suggesting potential speedups well beyond 2$\times$.

\subsection{Cost Savings at Scale}

If LST's 2$\times$ speedup transfers to frontier model training:

\begin{itemize}
    \item \textbf{Llama-3 405B:} Estimated training cost $\sim$\$100M $\rightarrow$ LST would save $\sim$\$50M.
    \item \textbf{Carbon footprint:} 50\% training time reduction translates directly to 50\% energy reduction.
    \item \textbf{Democratization:} Researchers with single-GPU setups can train models in half the wall-clock time.
\end{itemize}

\subsection{Limitations}

\begin{enumerate}
    \item \textbf{Single model scale:} We demonstrate LST on GPT-2 124M only. Scaling behavior to 1B+ parameter models requires validation.
    \item \textbf{Short training:} 2,000 steps is far from convergence. Acceptance rate trends suggest further decline at longer horizons.
    \item \textbf{Loss quality gap:} 12.5\% loss degradation exceeds our target of $<$1\%. Further hyperparameter tuning and the mitigations above (Section~X-B) may close this gap.
    \item \textbf{Single run:} Results are from one training run. Variance across seeds is not characterized.
    \item \textbf{Gradient checkpointing dependency:} The favorable cost ratio depends on gradient checkpointing making backward passes $\sim$3$\times$ more expensive than forward. Without checkpointing, the cost advantage diminishes.
\end{enumerate}

% ====================================================================
%  XII. CONCLUSION
% ====================================================================
\section{Conclusion}

We presented Learned Speculative Training (LST), a framework that replaces analytic weight extrapolation with an online-trained draft transformer for speculative training of LLMs. Our key contributions are:

\begin{enumerate}
    \item The learned draft model achieves 76.3\% acceptance rate on GPT-2 124M, substantially outperforming Leap+Verify's reported $\sim$9\% in chaotic regimes.
    \item Gradient accumulation amplifies the speculative cost advantage, creating an $O(A)$ standard-to-accepted cost ratio.
    \item Five engineering optimizations (batched BMM decode, zero CUDA syncs, eval-mode verify, snapshot-free rollback, and gradient accumulation) make the speculative overhead negligible.
    \item Measured 2.02$\times$ wall-clock speedup on A100 40GB (1,483.5s vs.\ 2,997.4s for 2,000 steps).
\end{enumerate}

Future work includes scaling to larger models, closing the loss quality gap via hybrid training schedules (Section~\ref{sec:ablations}), extending to multi-GPU distributed training where the speculative paradigm may yield additional communication savings, and characterizing the draft model's generalization across different model architectures and training corpora.

% ====================================================================
%  REFERENCES
% ====================================================================
\begin{thebibliography}{99}

\bibitem{kaplan2020scaling}
J.\ Kaplan, S.\ McCandlish, T.\ Henighan, T.\ P.\ Brown, B.\ Chess, R.\ Child, S.\ Gray, A.\ Radford, J.\ Wu, and D.\ Amodei, ``Scaling laws for neural language models,'' \textit{arXiv preprint arXiv:2001.08361}, 2020.

\bibitem{leviathan2023fast}
Y.\ Leviathan, M.\ Kalman, and Y.\ Matias, ``Fast inference from transformers via speculative decoding,'' in \textit{Proc.\ ICML}, 2023.

\bibitem{chen2023accelerating}
C.\ Chen, S.\ Borgeaud, G.\ Irving, J.-B.\ Lespiau, L.\ Sifre, and J.\ Jumper, ``Accelerating large language model decoding with speculative sampling,'' \textit{arXiv preprint arXiv:2302.01318}, 2023.

\bibitem{stern2018blockwise}
M.\ Stern, N.\ Shazeer, and J.\ Uszkoreit, ``Blockwise parallel decoding for deep autoregressive models,'' in \textit{Proc.\ NeurIPS}, 2018.

\bibitem{leapverify2026}
``Leap and Verify: Speculative training with learned weight predictions,'' \textit{arXiv preprint arXiv:2602.19580}, Feb.\ 2026.

\bibitem{metz2022velo}
L.\ Metz, J.\ Harrison, C.\ D.\ Freeman, A.\ Merchant, L.\ Beyer, J.\ Gilmer, and J.\ Sohl-Dickstein, ``VeLO: Training versatile learned optimizers by scaling up,'' \textit{arXiv preprint arXiv:2211.09760}, 2022.

\bibitem{jaderberg2017decoupled}
M.\ Jaderberg, W.\ M.\ Czarnecki, S.\ Osindero, O.\ Vinyals, A.\ Graves, D.\ Silver, and K.\ Kavukcuoglu, ``Decoupled neural interfaces using synthetic gradients,'' in \textit{Proc.\ ICML}, 2017.

\bibitem{baydin2022gradients}
A.\ G.\ Baydin, B.\ A.\ Pearlmutter, D.\ Syme, F.\ Wood, and P.\ Torr, ``Gradients without backpropagation,'' \textit{arXiv preprint arXiv:2202.08587}, 2022.

\bibitem{silver2025fwdgrad}
D.\ Silver et al., ``Forward gradient-based methods for efficient training,'' \textit{arXiv preprint}, 2025.

\bibitem{radford2019language}
A.\ Radford, J.\ Wu, R.\ Child, D.\ Luan, D.\ Amodei, and I.\ Sutskever, ``Language models are unsupervised multitask learners,'' \textit{OpenAI Blog}, 2019.

\bibitem{johnson1984extensions}
W.\ B.\ Johnson and J.\ Lindenstrauss, ``Extensions of Lipschitz mappings into a Hilbert space,'' \textit{Contemp.\ Math.}, vol.\ 26, pp.\ 189--206, 1984.

\bibitem{touvron2023llama}
H.\ Touvron et al., ``LLaMA: Open and efficient foundation language models,'' \textit{arXiv preprint arXiv:2302.13971}, 2023.

\end{thebibliography}

\end{document}
